% !PRINTED 225
\weberpage{225}
Hierzu ist noch zu bemerken, dass die überzählig eingeführten
Variabeln $t_n, t_{n+1} \ldots$ in der zweiten Form von $E_{0,s}, E_{1,s} \ldots E_{s-1,s}$
bereits wieder eliminirt sind, und dass die Variable $t_{n-1}$ nur
in $E_{s,s}$ vorkommt.

\begin{center}
\S.\,73.

\gesperrt{Die Bezoutiante.}
\end{center}

Die Gleichung (4) des vorigen Paragraphen wird, entwickelt,
die Gestalt haben
\[
y^n + P_1 y^{n-1} + P_2 y^{n-2} + \cdots + P_n = 0,
\]
oder wenn wir $t_{n-1}$ so bestimmen, dass $S(y) = 0$ wird,
\weq{1}{y^n + P_2 y^{n-2} + \cdots + P_n = 0;}
darin ist dann
\[
- 2 P_2 = S(y^2)
\]
eine Function zweiten Grades in Bezug auf die $t$ und in Bezug
auf die $a$. Diese Function ist für viele Anwendungen besonders
wichtig und hat von \gesperrt{Sylvester} den Namen Bezoutiante
der Function $f(x)$ erhalten, zu Ehren des französischen Mathematikers
\gesperrt{Bezout}, der schon im vorigen Jahrhundert die ersten
richtigen Ausführungen über Elimination gegeben hat.

Nach den Formeln des letzten Paragraphen lässt sich diese
Function verhältnissmässig einfach berechnen. Wir führen zunächst
neben den Variablen $t_{n-1}, t_{n-2} \ldots t_0$ noch ein zweites
System davon unabhängiger Variabeln $\tau_{n-1}, \tau_{n-2} \ldots \tau_0$ ein und
setzen
\weq{2}{
\begin{aligned}
y &= t_{n-1} f_0 + t_{n-2} f_1 + \cdots + t_0 f_{n-1} \\
z &= \tau_{n-1} f_0 + \tau_{n-2} f_1 + \cdots + \tau_0 f_{n-1}.
\end{aligned}
}
Statt nun $S(y^2)$ zu bilden, berechnen wir zunächst $S(yz)$,
woraus dann $S(y^2)$ hervorgeht, wenn man die $\tau$ gleich den $t$
setzt.

Wenn man die Formel (3) des vorigen Paragraphen
\weq{3}{y f_s = E_{0,s} f_0 + E_{1,s} f_1 + \cdots + E_{n-1,s} f_{n-1}}
mit $\tau_{n-s-1}$ multiplicirt und dann in Bezug auf $s$ von $0$ bis $n-1$
summirt, so findet man
\weq{4}{
\begin{aligned}
y z = {} & f_0 \sum_{0,\,n-1}^{s} E_{0,s} \tau_{n-s-1} \\
& + f_1 \sum_{0,\,n-1}^{s} E_{1,s} \tau_{n-s-1} + \cdots + f_{n-1} \sum_{0,\,n-1}^{s} E_{n-1,s} \tau_{n-s-1}.
\end{aligned}
}
